% !TEX root = ../main.tex
\chapter{Yöntem}
\label{chap:yontem}

Hata ölçütü olarak \ac{MPJPE} kullanılmıştır (\cref{eq:mpjpe}):
\begin{equation}
  \mathrm{MPJPE} = \frac{1}{N} \sum_{i=1}^{N}
  \left\| \hat{\mathbf{p}}_i - \mathbf{p}_i \right\|_2
  \label{eq:mpjpe}
\end{equation}

\begin{theorem}
  \label{thm:ornek}
  Her $x \in \mathbb{R}$ için $x^2 \geq 0$ sağlanır.
\end{theorem}

\begin{proof}
  Gerçel sayıların karesi negatif olamaz.
\end{proof}

\Cref{thm:ornek} ve \cref{fig:akis} yöntemin temelini oluşturur.

\begin{figure}[htbp]
  \centering
  \begin{subfigure}[b]{0.45\textwidth}
    \centering
    \includegraphics[width=\textwidth]{example-image-a}
    \caption{Girdi.}
    \label{fig:akis-girdi}
  \end{subfigure}
  \hfill
  \begin{subfigure}[b]{0.45\textwidth}
    \centering
    \includegraphics[width=\textwidth]{example-image-b}
    \caption{Çıktı.}
    \label{fig:akis-cikti}
  \end{subfigure}
  \caption{Yöntemin akışı (yer tutucu görseller).}
  \label{fig:akis}
\end{figure}
